problem:  
Let \( (M,g) \) be a compact Riemann surface and \( (N,h) \) a compact Riemannian manifold whose sectional curvature is nonpositive. Denote by  
\[
\partial_t u = \tau(u)
\]
the harmonic map flow \( u: M\times [0,T)\to N \) with smooth initial data \(u_0\).  Let \( \mathcal{M}(N) \) be the moduli space of harmonic two–spheres \( \phi:S^2\to N \), modulo reparametrization by \( \mathrm{PSL}_2(\mathbb{C}) \).  
For generic \(h\), one expects \(\mathcal{M}(N)\) to be a smooth manifold whose points are nondegenerate critical points of the energy restricted to \( \pi_2(N) \).

Suppose \( u_t \) develops a bubble tree in finite or infinite time.  The limit consists of a weakly harmonic map \(u_\infty:M\to N\) and finitely many harmonic spheres \( \{ \phi_i\}_{i=1}^k\subset \mathcal{M}(N)\) attached at distinct points of \(M\).  

**Question:**  
For a residual (Baire generic) subset of metrics \(h\) on \(N\), does the following rigidity principle hold?

1. *(Analytic rigidity)* Each $\phi_i$ appearing as a bubble must be a **nondegenerate** critical point of the energy (i.e., its Jacobi operator on normal variations has no zero modes).  
2. *(Homotopical rigidity)* The total homotopy class \([u_0]\in [M,N]\) determines, up to finite ambiguity, the collection of homotopy classes \([\phi_i]\in \pi_2(N)\) that can appear as bubbles, together with their multiplicities.  
3. *(Moduli rigidity)* All bubble trees corresponding to the same \([u_0]\) and generic \(h\) are equivalent in the sense that they correspond to the same set of nondegenerate points in \(\mathcal{M}(N)\), possibly differing only by reparametrization of the domain spheres.  

Equivalently, does there exist for generic \(h\) a *finite stratification* of \([M,N]\) by bubble–tree types determined by nondegenerate harmonic spheres in \(N\)?  

Why is it a "good" problem:  
This problem formalizes the idea that bubbling in harmonic map flow might be completely “discretized” for generic target geometries—each energy concentration corresponds to a nondegenerate harmonic sphere, with no moduli freedom, and the resulting bubble–tree type depends only on homotopy‐theoretic input. It merges analytic aspects (regularity and transversality of the harmonic map equation), differential–topological aspects (structure of \( [M,N] \) and \( \pi_2(N) \)), and geometric measure–theoretic aspects (compactification of moduli by bubble trees). Establishing such a rigidity would unify Morse theory on mapping spaces with geometric analysis of singular PDE flows, opening a bridge between variational topology and global analysis. The statement is simple and natural, yet remains far beyond current techniques, satisfying the criteria of depth, cross–disciplinary connection, and conceptual simplicity that define a truly “good” mathematical problem.